\documentclass{beamer}
% Setup authorship. This is placed in the preamble because it is not supposed
% to change much...
% NOTA BENE: Use the \inst{?} command only if the authors have different
% affiliation.
\author
{Óscar Pereira\inst{1,2} \and John Doe\inst{2}\\[0.75em]
  \texttt{\scriptsize \{https://, oscar@\}randomwalk.eu}\\
  \texttt{\scriptsize \{https://, john@\}somedoe.com} 
}
\institute
{
  \inst{1}%
  Some Institution
  \and
  \inst{2}%
  Some other Institution
}

% Fonts.
\usepackage{times}
\usepackage{fontspec}
\defaultfontfeatures{Ligatures=TeX}

% Language: replace UKenglish with portuguese or other lang.
\usepackage[UKenglish]{babel}

\usepackage{sfmath} % Math in sans serif font.
\usepackage{tikz} % To draw pictures!
\usepackage{textpos} % To draw in absolute positions.
\usepackage{marvosym} % For \Smiley, and LOTS of other things

% In slides, emphasis means italic AND bold face.
\renewcommand{\emph}[1]{\textbf{\textit{#1}}}

% For bibliography: in presentations, [author, year] is more meaningful then
% [number].
\usepackage[authoryear]{natbib} 

% Gets rid of bottom navigation symbols.
\setbeamertemplate{navigation symbols}{}

% Show page numbers.
\setbeamertemplate{footline}[frame number]

% Setup theme, colours, and transparency use.
\mode<presentation>
{
  \usetheme{default}
  \usecolortheme{whale}

  % for transparencies
  \setbeamercovered{transparent}
}

% If you have a file called "university-logo-filename.xxx", where xxx
% is a graphic format that can be processed by latex or pdflatex,
% resp., then you can add a logo as follows:

% \pgfdeclareimage[height=0.5cm]{university-logo}{university-logo-filename}
% \logo{\pgfuseimage{university-logo}}

% -------------------------------------------------------
% -------------------------------------------------------
% ------------------- CUSTOM STUFF ----------------------
% -------------------------------------------------------
% -------------------------------------------------------

%
% Packages.
%

%
% (Re)new commands. (Except those that are part of a specific setup).
%
\newcommand{\emd}{\textemdash}

% For § symbol. Some TeX compilers might complain if it is entered literally,
% so use this command instead.
\newcommand{\ts}{\textsection}

%
% All other custom stuff.
%

% -------------------------------------------------------
% -------------------------------------------------------
% ------------------- END CUSTOM STUFF ------------------
% -------------------------------------------------------
% -------------------------------------------------------


\title
{A Presentation Template}

\subtitle
{With a Subtitle, if Needed}

\date
{March 2020}

\begin{document}

\begin{frame}[plain]
  \titlepage
\end{frame}

\begin{frame}{On Science}
  \vfill
  \begin{itemize}
    \item<1-> Science is \emph{cool!} \citep{Sagan}
      \vfill
    \item<2-> Unfortunately \emph{a big great many lot} of people do not seem to get it\dots \citep{Feynman74}
      \vfill
    \item<3-> And now some \emph{math:}
      \vfill
      \begin{theorem}[Euler's identity] % Note the ( and ) in the optional name
        \begin{equation}
          e^{i\pi}+1=0
        \end{equation}
      \end{theorem}
  \end{itemize}
  \vfill
\end{frame}

\begin{frame}[fragile]{Verbathings}
  \begin{itemize}
    \item Not verbatim text, vs.\ \verb+verbatim text+
    \item The \verb+verbatim+ \emph{environment} also works:
      \begin{verbatim}
Lorem ipsum dolor sit amet, blah blah...
      \end{verbatim}
  \end{itemize}
\end{frame}

\begin{frame}{Summary}
  \begin{itemize}
  \item
    The \alert{first main message} of your talk in one or two lines.
  \item
    The \alert{second main message} of your talk in one or two lines.
  \item
    Perhaps a \alert{third message}, but not more than that.
  \end{itemize}
\end{frame}

\begin{frame}{References}
  \vfill
  \bibliography{sources}
  \bibliographystyle{plainnat}
  \vfill
\end{frame}

\begin{frame}[plain]{Thank you for listening!}
  \begin{textblock*}{0.1\textwidth}(0.2\textwidth, 0.25\textheight)
    \LARGE \textit{Questions}
  \end{textblock*}
  \begin{centering}
    \hfill\hfill
		\begin{tikzpicture}[ultra thick,scale=0.7]
			\draw[line width=3.5,line cap=round] (1.5,0) .. controls ++(0,4) and 
			++(0,-2) .. (4,4)
			to[out=90,in=0] (2,6)
			to[out=180,in=90] 
			(0,4);
    \end{tikzpicture}
    \hfill
  \end{centering}

  \begin{textblock*}{0.1\textwidth}(0.55\textwidth, 0.02\textheight)
    \Huge \Smiley
  \end{textblock*}
\end{frame}

\end{document}
